\section{Data: Policy Menus from the Bluebooks}\label{sec:data}

\subsection{How the FOMC prepares a decision}

The Federal Open Market Committee (FOMC) meets eight times a year at scheduled meetings and occasionally by
conference call. Each scheduled meeting is preceded by a structured set of staff documents
\citep{danker2005}. About two weeks before the meeting the \emph{Beige Book} summarises anecdotal evidence from
the twelve Reserve Bank districts. About a week before, the \emph{Greenbook} gives the Board staff's forecast for
the domestic and international economy. On the Friday before the meeting the \emph{Bluebook} sets out the policy
alternatives for the Committee's consideration, usually labelled A, B and C, each with a proposed federal funds
rate target and a rationale. The meeting agenda then includes a discussion of these alternatives before the
Committee votes. The format of these documents has changed over time: documents have been introduced, merged and
retired since the corpus begins in 1936. Over the sample used here the sequence above is stable.

Two of these documents enter the analysis. The Bluebook defines the \emph{menu} of actions the Committee
considered, which is the object of this thesis, and the Greenbook provides the staff forecasts used as
covariates in Section~\ref{sec:estimation}.

\subsection{Extracting the menus}

The Bluebooks are published with a five-year lag on the Board's website as HTML and PDF files. Building on the
FOMC text-analysis project of Montiel Olea, Giesecke, Chitale and Frermann, the documents were converted to raw
text and every sentence mentioning a policy alternative was extracted with the regular expression
\begin{center}
\small\verb/([^.]*)(alternatives?\s+[a-e])(\.|[^a-z]\.|[^a-z][^\.]+\.)/
\end{center}
which captures a sentence containing ``alternative'' followed by a letter. The extracted sentences were then
read and classified by hand, because the same target is sometimes described in different ways under different
rationales. For each alternative I record the proposed change in the federal funds rate target, and classify it
as easing ($E$), unchanged ($U$) or tightening ($T$).

The Committee's actual decision is taken from the target changes in the replication files of
\citet{angrist2018}. The sample runs from August 1989 to January 2006, the Greenspan chairmanship after the
Committee had begun to target the federal funds rate explicitly \citep{romer2004}. It contains 133 scheduled
meetings with a Bluebook menu. Target changes of 25, 50 and 75 basis points are pooled into the three categories
because 50 and 75 basis point moves are too rare to estimate separately.

\subsection{What the menus show}

Table~\ref{tab:menus} cross-tabulates the menu offered with the decision taken. Two facts stand out. First, the
decision is \emph{always} on the menu: no meeting in the sample took an action that the Bluebook did not propose.
The zeros in Table~\ref{tab:menus} are therefore structural, and they are observed rather than inferred. Second, most
menus are restricted: 87 of 133 meetings (65\%) were offered a menu without easing or without tightening, and only
46 had all three options. The full menu is also where the Committee was least likely to move. Only one of the 27
tightenings in the sample came from a full menu.

\begin{table}[!htbp]
  \centering
  \caption{Bluebook menus and FOMC decisions, August 1989 -- January 2006}\label{tab:menus}
  \tablefile{menus.tex}
  \par\medskip
  \begin{minipage}{0.8\linewidth}\footnotesize
  \emph{Notes:} Scheduled FOMC meetings. Rows give the set of alternatives proposed in the Bluebook ($E$ ease,
  $U$ unchanged, $T$ tighten); columns give the decision. Every decision lies in the menu.
  \end{minipage}
\end{table}

Figure~\ref{fig:state} shows where in the macroeconomic state space these restrictions bind. Panel (a) plots
lagged inflation and industrial-production growth at each meeting, marking the meetings at which the Committee
tightened or eased. Panel (b) marks instead the meetings whose menu \emph{excluded} easing or tightening. Menus
without easing are concentrated where output growth is above its sample average (88\% of them), and menus without
tightening where it is below (78\%). The same pattern appears in the decisions themselves, but the menu shows it
without relying on a fitted model: in about half of the meetings with below-average output growth (29 of 56) tightening
was not even proposed, and the Committee tightened at only three of them. This is
exactly the kind of covariate region in which a propensity score must be zero, and it is why a standard IPW
estimator, which divides by the propensity score, is not available there.

\begin{figure}[!htbp]
  \centering
  \includegraphics[width=\linewidth]{state_space.pdf}
  \caption{Policy decisions and policy menus in the inflation--output space. Each point is a scheduled FOMC
  meeting; axes show year-on-year PCE inflation and industrial-production growth in the previous month, as
  deviations from their 1989--2006 means. Grey points are meetings not in the highlighted group.}\label{fig:state}
\end{figure}
